\documentclass{article}
\usepackage[T1]{fontenc}
\usepackage{lmodern}
\usepackage{graphicx}
\usepackage{amsmath,amssymb}
\usepackage[a4paper,margin=2cm]{geometry}
\usepackage[absolute,overlay]{textpos}
\setlength{\TPHorizModule}{1mm}
\setlength{\TPVertModule}{1mm}
\usepackage[indonesian]{babel}
\usepackage{hyperref}
\setlength{\emergencystretch}{2em}
% unit: Halaman 58

\begin{document}

% === Halaman 58 (python/native) ===

Tiga (3) Jenis Algoritma Kriptografi

1. Algoritma kriptografi simetri (symmetric-key cryptography)

•  Kunci enkripsi = kunci dekripsi, harus dijaga privat (rahasia) atau secret

•   Sudah ada sejak ribuan tahun yang hingga tahun 1976   

58

Kunci privat, K

Enkripsi

EK (P) = C

Dekripsi

DK (C) = P

Cipherteks, C

Plainteks, P

Plainteks, P

Kunci privat, K

-

Data Encryption Standard (DES)

-

Advanced Encryption Standard (AES)

-

Serpent

-

Blowfish

-

Loki

-

MARS

-

RC6

-

Twofish

-

3-DES

-

IDEA

-

FEAL

-

RC4

-

SEAL

-

Panama

-

dll

\end{document}
